\thispagestyle{otFront}
\phantomsection\pdfbookmark[0]{Parcours de travail}{bm-parcours}
\otPageTitre{Parcours de travail}{Des bases aux problèmes de recherche}
\fontsize{11}{15}\selectfont
\otIntertitre{Un livre de Première générale, spécialité mathématiques}
Les chapitres 1 à 11 développent l'algèbre, l'analyse, la géométrie, les probabilités
et les statistiques. Les chapitres 12 à 14 donnent les outils transversaux : Python,
raisonnement et automatismes. Ces derniers sont à reprendre au fil de l'année.
L'ordre des thèmes peut être adapté à la progression du professeur.\par\vspace{5mm}
\otIntertitre{Choisir le bon point de départ}
\begin{tabularx}{\linewidth}{@{}p{38mm}Y@{}}\toprule
\textbf{Parcours} & \textbf{Objectif}\\\midrule
Découverte, 1 à 5 & Vérifier les bases avec des applications directes.\\[2mm]
Entraînement, 6 à 11 & Installer les méthodes, puis choisir une démarche.\\[2mm]
Problèmes, 12 à 15 & Relier les notions, modéliser et rédiger.\\[2mm]
Défis D01 à D14 & Prendre des initiatives et traiter des discussions exigeantes.\\[2mm]
Sujets blancs A et B & Travailler en temps limité, sans calculatrice.\\\bottomrule
\end{tabularx}
\vspace{6mm}
\otIntertitre{Une méthode de travail en quatre temps}
\begin{questions}
\item Lire les objectifs et le cours. Reformuler une propriété avec ses hypothèses.
\item Chercher sans corrigé. Écrire les données, un schéma utile et les premières relations.
\item Comparer la démarche à la correction. Repérer précisément l'étape qui manquait.
\item Refaire quelques jours plus tard, puis expliquer la solution à voix haute.
\end{questions}
\vspace{4mm}
\otIntertitre{Se repérer et mesurer ses progrès}
Le tableau de bord suit les 210 exercices des chapitres. Pour les 14 défis,
noter dans son cahier : réussi seul, réussi avec un indice, ou à reprendre.
Les indications sont séparées des énoncés et des corrigés.
Les démonstrations de référence et les repères du programme servent de points d'appui.
\editionRetenir{La difficulté augmente par les choix de méthode et le raisonnement.
Un défi facultatif peut demander du temps, tout en restant fondé sur les outils de Première.}
\clearpage
